\section{Integral, irreducible, and reduced schemes}
\label{section-integral}

\begin{definition}
\label{definition-integral}
Let $X$ be a scheme. We say $X$ is {\it integral} if it is nonempty and
for every nonempty affine open $\Spec(R) = U \subset X$ the ring $R$
is an integral domain.
\end{definition}

\begin{lemma}
\label{lemma-characterize-reduced}
Let $X$ be a scheme.
The following are equivalent.
\begin{enumerate}
\item The scheme $X$ is reduced, see
Schemes, Definition \ref{schemes-definition-reduced}.
\item There exists an affine open covering $X = \bigcup U_i$
such that each $\Gamma(U_i, \mathcal{O}_X)$ is reduced.
\item For every affine open $U \subset X$ the ring
$\mathcal{O}_X(U)$ is reduced.
\item For every open $U \subset X$ the ring $\mathcal{O}_X(U)$ is reduced.
\end{enumerate}
\end{lemma}

\begin{proof}
See Schemes, Lemmas \ref{schemes-lemma-reduced} and
\ref{schemes-lemma-affine-reduced}.
\end{proof}

\begin{lemma}
\label{lemma-characterize-irreducible}
Let $X$ be a scheme.
The following are equivalent.
\begin{enumerate}
\item The scheme $X$ is irreducible.
\item There exists an affine open covering $X = \bigcup_{i \in I} U_i$
such that $I$ is not empty, $U_i$ is irreducible for all $i \in I$, and
$U_i \cap U_j \not = \emptyset$ for all $i, j \in I$.
\item The scheme $X$ is nonempty and every nonempty affine open
$U \subset X$ is irreducible.
\end{enumerate}
\end{lemma}

\begin{proof}
Assume (1). By Schemes, Lemma \ref{schemes-lemma-scheme-sober}
we see that $X$ has a unique generic point $\eta$. Then
$X = \overline{\{\eta\}}$. Hence $\eta$ is an element of
every nonempty affine open $U \subset X$. This implies
that $\eta \in U$ is dense hence $U$ is irreducible.
It also implies any two nonempty affines meet.
Thus (1) implies both (2) and (3).

\medskip\noindent
Assume (2). Suppose $X = Z_1 \cup Z_2$ is a union of two closed subsets.
For every $i$ we see that either $U_i \subset Z_1$ or $U_i \subset Z_2$.
Pick some $i \in I$ and assume $U_i \subset Z_1$ (possibly after renumbering
$Z_1$, $Z_2$). For any $j \in I$ the open subset $U_i \cap U_j$ is dense in
$U_j$ and contained in the closed subset $Z_1 \cap U_j$. We conclude that
also $U_j \subset Z_1$. Thus $X = Z_1$ as desired.

\medskip\noindent
Assume (3). Choose an affine open covering $X = \bigcup_{i \in I} U_i$.
We may assume that each $U_i$ is nonempty.
Since $X$ is nonempty we see that $I$ is not empty.
By assumption each $U_i$ is irreducible.
Suppose $U_i \cap U_j = \emptyset$ for some pair $i, j \in I$.
Then the open $U_i \amalg  U_j = U_i \cup U_j$ is affine, see
Schemes, Lemma \ref{schemes-lemma-disjoint-union-affines}.
Hence it is irreducible by assumption which is absurd. We conclude that (3)
implies (2). The lemma is proved.
\end{proof}

\begin{lemma}
\label{lemma-characterize-integral}
A scheme $X$ is integral if and only if it is reduced and irreducible.
\end{lemma}

\begin{proof}
If $X$ is irreducible, then every affine open $\Spec(R) = U \subset X$
is irreducible. If $X$ is reduced, then $R$ is reduced, by
Lemma \ref{lemma-characterize-reduced} above. Hence $R$ is reduced
and $(0)$ is a prime ideal, i.e., $R$ is an integral domain.

\medskip\noindent
If $X$ is integral, then for every nonempty affine open
$\Spec(R) = U \subset X$ the ring $R$ is reduced
and hence $X$ is reduced by Lemma \ref{lemma-characterize-reduced}.
Moreover, every nonempty affine open is irreducible.
Hence $X$ is irreducible, see Lemma \ref{lemma-characterize-irreducible}.
\end{proof}

\noindent
In Examples, Section
\ref{examples-section-connected-locally-integral-not-integral}
we construct a connected affine scheme all of whose local rings are domains,
but which is not integral.











